\section{Discussion}
\label{sec:discussion}
\subsection{Complementary local coverage}
The principal benefit of the heterogeneous local branch is complementary coverage. Under the final OR rule, locally detected attacks improve the combined decision only when the supervised peer union misses them. This explains why the gains depend on the availability of peer-missed attacks and why improved standalone local detection need not translate into proportional end-to-end gains. Local-exclusive coverage at $\alpha=0.01$ and $N=400$ reaches 14.74\% in TON\_IoT and 12.10\% in IoT-23, whereas it is nearly zero in N-BaIoT.

The rarity-family removals produce the largest detection losses among the tested expert-family ablations in all four primary configurations. This supports the usefulness of marginal and pairwise feature-rarity evidence alongside the shape- and neighborhood-based detectors. The smaller response-expert effects indicate that joint peer-response modeling supplies additional, but less extensive, coverage. Because each removal recalibrates the remaining union, the losses characterize complete ablated rules rather than an additive attribution to individual experts.

The protected anchor preserves an averaged local decision path alongside sparse expert alerts. Its value is an operating tradeoff: at $\alpha=0.05$, removing it raises both detection and false alerts. The primary comparison also changes the comparator's seven-member composition and thresholding, so the overall gain belongs to the complete local-branch replacement.

\subsection{Retaining peer evidence separately}
Peer-wise calibration gives each classifier its own target-relative threshold. Unlike a threshold on a fused maximum, a high threshold for one peer need not directly exclude another peer's smaller evidence. This addresses the peer-drowning mechanism while retaining dependence through the common union multiplier.

The peer-fusion ablations retain the PACE local bank and test this architectural choice separately from the primary local-branch comparison. At $\alpha=0.01$, PACE exceeds peer-mean fusion by 7.34 and 5.60 percentage points and peer-maximum fusion by 4.56 and 3.64 points at $N=100$ and $400$. At $\alpha=0.05$, the corresponding advantages are 5.25 and 4.47 points over the mean, and 3.05 and 3.16 over the maximum. All eight paired intervals exclude zero. The fused arms also have lower mean FPR: 0.023855--0.026093 at $\alpha=0.01$ and 0.043049--0.047157 at $\alpha=0.05$, compared with PACE's 0.032369--0.033992 and 0.054337--0.055453. These results support a detection advantage at the tested nominal parameters, accompanied by an alert-rate tradeoff.

The shared multiplier uses overlap between calibration alerts to relax thresholds while checking their union on every block. It does not require independent peer alerts. It also couples the retained thresholds, since changing one expert's calibration alerts can change the multiplier for all peers. The empirical evidence transformation and union search retain the calibration interpretation specified in Section~\ref{sec:method}.

\subsection{Variation across targets and fitting budgets}
TON\_IoT offers the largest observed gain and substantial local-exclusive coverage. IoT-23 also benefits on average, although its gains are concentrated in a subset of targets. N-BaIoT's near-complete peer coverage leaves little room for the local path to improve the final union. CTU-13 has small and uncertain stratum gains and includes the scenario-12 regression. These outcomes locate the empirical value of complementarity without attributing corpus differences to unmeasured traffic properties.

Increasing local fitting data changes the learned local evidence, but the final detector benefits only from changes outside existing peer coverage. On the fixed 58-target cohort, moving from 100 to 400 rows produces small, uncertain end-to-end changes; the 750-row comparison still favors PACE over the comparator. Benefits are weaker in the 47-target and 15-target resource-eligible cohorts, and the high-resource $\alpha=0.05$ budget change is negative. Consequently, neither larger budgets nor higher standalone local coverage imply monotonic improvement of the final OR rule.

\subsection{Practical interpretation}
At $\alpha=0.05$, PACE's primary detection gains accompany lower mean FPR at both budgets. Operational assessment must additionally consider the target distribution: at $\alpha=0.01$, medians below 0.01 coexist with means above 0.03 and worst-target rates near 0.67. The mean improvement is useful comparative evidence, while the upper tail shows why the nominal parameter cannot serve as a uniform operational guarantee.

Applying the detector requires trusted benign onboarding data for fitting, reference construction, and calibration, plus the preprocessing access in Table~\ref{tab:data-access}. References and thresholds remain fixed in the evaluated rule. Independent validation would need to establish when to update them, how to respond to excessive alerts, and how any recalibration changes detection.

Row alerts also require a policy for aggregation into actionable incidents. An anomalous record need not represent an actionable security event~\cite{sommer2010closed}, and precision depends on attack prevalence~\cite{saito2015precision}. The benchmark's separately sampled classes and target-macro rates do not measure incident precision or analyst workload. These operational questions remain distinct from the demonstrated detection contribution.
